\section{Priority research goal: robust finite-time cooling}
\label{sec:physicalfinitetimegoal}
\textbf{Research priority recorded 12 October 2026; verified rate-model outcome below.}
The objective is a physically motivated thermalization rate model with a
substantial, robust \emph{finite-time} Mpemba advantage under moderate parameter
ranges. This takes priority over extending the asymptotic amplitude extremum.
The sharp $25\%$ bound concerns a slow-mode amplitude; it does not establish
a $25\%$ reduction in cooling time. The existing extreme-parameter sharpness
families are therefore mathematical benchmarks, not solutions of this goal.

\subsection{Physical model and scope}
Start with thermally activated transitions between four metastable states,
using the energy-landscape framework of
\href{https://pmc.ncbi.nlm.nih.gov/articles/PMC5441807/}{Lu and Raz (2017)}.
For a transition from $j$ to $i$, use
\[
 a_{ij}=\nu_{ij}\exp[-\beta_b(B_{ij}-E_j)],\qquad
 B_{ij}=B_{ji}>\max(E_i,E_j),\qquad \nu_{ij}=\nu_{ji}>0.
\]
Then $a_{ij}e^{-\beta_b E_j}=a_{ji}e^{-\beta_b E_i}$, so detailed balance
is automatic. Both preparations evolve with this same bath generator.
Begin with a common attempt frequency and a low-dimensional, physically
specified barrier geometry. Choose that geometry before optimizing the
effect. Merely expressing an arbitrary reversible matrix as six independently
fitted barriers does not establish a physical mechanism. A convincing
candidate must explain its state interpretation, barrier constraints,
metastable coarse-graining and separation from intrawell relaxation; a
continuum landscape or other independent microscopic model should test the
rate approximation when feasible.

The preferred extension of the present paper also retains positive rates,
a simple slow mode and failure of every original-rate triplet test.
Check and report those predicates separately. A physically useful example
that fails this additional premise may meet the finite-time research
objective, but is not a stronger counterexample to the triplet conjecture.
Similarly, sparse topologies and additional state counts must be labeled as
changes of scope, not silently substituted for the four-state construction.

\subsection{Provisional quantitative benchmarks}
These are initial research targets, not established experimental standards.
Set $k_B T_b=1$ and fix the reference attempt frequency $\nu_0$.
Use ordered energies with adjacent gaps at least $0.1$ and total span at most
$5$; take barrier excesses $B_{ij}-\max(E_i,E_j)$ between $2$ and $6$.
Start with $\nu_{ij}=\nu_0$, and require the ratio of the largest to smallest
nonzero directed rate to be at most $10^3$. Restrict preparations to
$1.25\leq T_c/T_b<T_h/T_b\leq10$.
These constraints jointly define the initial search box; they do not
guarantee that any target is feasible. Report the actual rate range and
relaxation time, not just the energy range. Any wider box, variable
prefactors, or changed threshold must be recorded as a separate study.

For $x\in\{h,c\}$, let $p_x(t)=e^{Mt}\pi(T_x)$ and use relative entropy
in nats as the primary distance,
\[
 D_x(t)=\sum_i p_{x,i}(t)\log\frac{p_{x,i}(t)}{\pi_{b,i}},
 \qquad
 A_D(t)=1-\frac{D_h(t)}{D_c(t)},\qquad
 \tau=\frac{1}{|\lambda_2|}.
\]
Require normal initial ordering $D_h(0)>D_c(0)$. Target both:
\begin{enumerate}
\item \textbf{A resolved distance advantage.}
Find a contiguous interval $[t_1,t_2]\subset(0,5\tau]$ of length at least
$0.25\tau$ on which $A_D(t)\geq0.10$ and $D_c(t)\geq10^{-3}$.
The distance floor excludes a large percentage arising only in an
arbitrarily small asymptotic tail.
\item \textbf{Time saved to a common threshold.}
Define $t_x(\delta)=\inf\{t\geq0:D_x(t)\leq\delta\}$ and
$A_t(\delta)=1-t_h(\delta)/t_c(\delta)$.
Target $A_t(\delta)\geq0.05$ throughout the predeclared band
$10^{-3}\leq\delta\leq10^{-2}$, with both initial distances above $10^{-2}$.
Report absolute time savings in units of $\nu_0^{-1}$ as well as percentages
and $\tau$ units. These are post-quench relaxation times; preparation time
and energy cost require a separate accounting for any operational claim.
\end{enumerate}
Also report total-variation and Pearson chi-square trajectories to expose
distance dependence. A KL success is not automatically a success for every
distance. Common rate rescaling cannot improve either percentage and must
not be presented as a cooling mechanism.

\subsection{Robustness, verification and decision record}
The target is a neighborhood, not a single tuned optimum. Test simultaneous
$\pm5\%$ variations in the independent positive energy gaps, barrier
excesses and preparation-temperature excesses $T_x/T_b-1$, rebuilding
detailed balance each time. For models with independently variable
prefactors, perturb those too. Specify parameter correlations imposed by
the physical geometry. Seek a neighborhood inside the declared parameter
box in which both performance benchmarks survive; quantify a smaller
verified neighborhood if the target tolerance fails. If the result is
advertised as an extension of the counterexample, the triplet predicates
must also survive throughout that neighborhood.

Use sampled perturbations and optimization for discovery, followed by
independent high-precision matrix-exponential and first-passage checks.
Certify interval-wide distance and threshold claims with interval bounds
or analytic estimates; a time grid or a random perturbation cloud alone
does not certify robustness. Record seeds, search boxes, tolerances,
rejected candidates and all benchmark revisions. The deliverable is a
compact physical model, explicit parameters, a finite-time performance
figure, reproducible calculations and a stated robustness domain, with
physical approximation errors distinguished from numerical errors.

The first milestone was to specify and justify the constrained barrier
geometry, then screen the declared box and analyze the best candidates.
Sections~\ref{sec:physicalresults}--\ref{sec:thermalcollision} now establish
both benchmarks on the full requested neighborhood for an idealized
thermal-collision rate model. Material and continuum validation remain
outside that verified conclusion.
At the time this priority was recorded, no new candidate meeting these
benchmarks had been established. If a model or box fails, report the
bounded search evidence or a proved obstruction and the next justified
model change. Failure to find a candidate is not a general impossibility
theorem. New evidence and any resulting conclusions belong in this
comprehensive research report before being promoted into the focused paper.

